\chapter{Long QT syndrome}
\index{Long QT syndrome}

\section*{Definition and Overview}
Long QT syndrome (LQTS) is a rare channelopathy characterised by a prolonged QT interval on the electrocardiogram (ECG) and an increased risk of torsades de pointes, a life-threatening polymorphic ventricular tachycardia. This prolongation of the ventricular action potential is caused by dysfunction in cardiac ion channels, which can be congenital (inherited) or acquired (secondary to medications or electrolyte disturbances). LQTS can lead to sudden, unexplained episodes of syncope, seizures, or cardiac arrest, often precipitated by physical exertion or emotional stress. Early diagnosis and individualised management are critical to preventing sudden cardiac death, particularly in young, otherwise healthy individuals.

\section*{Epidemiology}
Congenital LQTS has an estimated prevalence of approximately 1 in 2,000 live births worldwide, making it one of the more common inherited arrhythmia syndromes. It affects males and females, though symptoms typically present earlier in males (often before puberty) whereas females are at a higher risk of events during adolescence and the post-partum period. In many countries, LQTS is a recognised cause of sudden unexplained death in young people, including a proportion of cases historically attributed to sudden infant death syndrome (SIDS) or unexplained drowning. Acquired LQTS is far more common than the congenital form, frequently occurring in hospitalised patients exposed to QT-prolonging medications, with a higher incidence in elderly individuals and those with pre-existing structural heart disease.

\section*{Aetiology and Pathophysiology}
The cardiac action potential relies on a delicate balance of inward depolarising (sodium and calcium) and outward repolarising (potassium) currents. LQTS arises when this balance is disrupted, delaying ventricular repolarisation.
\begin{itemize}
    \item \textbf{Congenital genetic mutations:} Inherited mutations in genes encoding cardiac ion channel subunits. The three most common genotypes are LQT1 (\textit{KCNQ1} gene, accounting for 30\% to 35\% of cases), LQT2 (\textit{KCNH2} gene, 25\% to 30\%), and LQT3 (\textit{SCN5A} gene, 5\% to 10\%).
    \item \textbf{Impaired potassium efflux (LQT1 and LQT2):} Mutations in \textit{KCNQ1} and \textit{KCNH2} reduce the function of outward-rectifying potassium channels ($I_{Ks}$ and $I_{Kr}$ currents, respectively), slowing down the exit of potassium ions and prolonging repolarisation.
    \item \textbf{Excessive sodium influx (LQT3):} Mutations in \textit{SCN5A} cause a "gain-of-function" in the cardiac sodium channel ($Na_V1.5$), leading to a persistent inward sodium current ($I_{Na}$) during the plateau phase of the action potential, which delays repolarisation.
    \item \textbf{Early afterdepolarisations (EADs):} Prolonged action potentials allow L-type calcium channels to recover from inactivation and reopen, generating EADs. If these EADs reach threshold, they trigger premature ventricular contractions, which can initiate torsades de pointes.
    \item \textbf{Acquired triggers:} Medications (such as certain antiarrhythmics, antibiotics, antipsychotics, and antihistamines) that block the hERG potassium channel ($I_{Kr}$), or metabolic abnormalities (such as hypokalaemia, hypomagnesaemia, and hypocalcaemia).
\end{itemize}

\section*{Clinical Features}
Many individuals with LQTS are asymptomatic, but when clinical features occur, they are typically related to transient ventricular arrhythmias.
\begin{itemize}
    \item \textbf{Unexplained syncope:} Sudden fainting spells, often without warning or prodromal symptoms, typically triggered by exercise, emotional excitement, or sudden auditory stimuli (such as an alarm clock or telephone ring).
    \item \textbf{Seizure-like activity:} Transient cerebral hypoperfusion during prolonged episodes of torsades de pointes can mimic generalised tonic-clonic seizures, often leading to a misdiagnosis of epilepsy.
    \item \textbf{Cardiac arrest and sudden death:} Ventricular fibrillation developing from degenerated torsades de pointes, presenting as sudden cardiac arrest, which may be the first clinical manifestation of the syndrome.
    \item \textbf{Electrocardiographic abnormalities:} Prolonged rate-corrected QT interval (QTc typically $>470$~ms in females and $>450$~ms in males), bizarre T-wave morphology (notched, biphasic, or low amplitude), and prominent U waves.
    \item \textbf{Precipitating triggers by genotype:} LQT1 events are classically triggered by swimming or physical exertion; LQT2 events by sudden loud noises or emotional stress; and LQT3 events typically occur during sleep or rest when the heart rate is low.
    \item \textbf{Palpitations and dizziness:} Sensation of rapid, irregular heartbeats or lightheadedness, representing self-terminating runs of torsades de pointes.
\end{itemize}

\section*{Differential Diagnosis}
LQTS must be distinguished from other causes of syncope, seizures, and sudden cardiac arrest to ensure correct treatment and avoid inappropriate therapies.
\begin{itemize}
    \item \textbf{Vasovagal syncope:} The most common cause of fainting, typically preceded by a clear prodrome (nausea, sweating, warmth) and triggered by prolonged standing or pain, with a normal ECG.
    \item \textbf{Epilepsy:} Presents with recurrent seizures, but is associated with abnormal electroencephalogram (EEG) findings and post-ictal confusion, whereas LQTS-induced syncopal seizures have rapid recovery once perfusion is restored.
    \item \textbf{Catecholaminergic polymorphic ventricular tachycardia (CPVT):} Another inherited arrhythmia syndrome triggered by exercise or stress, but presents with a normal resting QTc interval and stress-induced bidirectional ventricular tachycardia.
    \item \textbf{Brugada syndrome:} An inherited channelopathy presenting with syncope or sudden death, but characterised by ST-segment elevation in the right precordial leads (V1-V3) rather than QTc prolongation.
    \item \textbf{Arrhythmogenic right ventricular cardiomyopathy (ARVC):} A structural heart disease causing ventricular arrhythmias, distinguished by fibrofatty replacement of the right ventricle on echocardiogram or cardiac MRI.
\end{itemize}

\section*{When to Seek Medical Care}
Immediate medical evaluation is required for anyone who experiences unexplained fainting, especially during exercise or emotional stress, or if there is a family history of sudden, unexplained death in young people.

\textbf{Contact your local emergency services for:}
\begin{itemize}
    \item Sudden loss of consciousness or fainting.
    \item A seizure, particularly if it occurs during or immediately after exercise or sudden excitement.
    \item Symptoms of cardiac arrest, where the person is unresponsive and not breathing normally, requiring immediate CPR and defibrillation.
\end{itemize}
Other non-emergency reasons to seek medical care include palpitations, lightheadedness, or when a family member has been diagnosed with LQTS, necessitating screening.

\section*{Diagnosis and Investigations}
The diagnosis of LQTS involves a combination of clinical scoring, electrocardiographic evaluation, and genetic testing, using the Schwartz score criteria.
\begin{itemize}
    \item \textbf{12-lead electrocardiogram (ECG):} The primary diagnostic tool to measure the QT interval, corrected for heart rate using Bazett's formula ($QTc = QT/\sqrt{RR}$), with repeated tracings often needed due to QTc variability.
    \item \textbf{Exercise stress testing:} Useful in LQT1, where the QTc interval typically fails to shorten or actually prolongs during exercise and recovery, helping to unmask the diagnosis.
    \item \textbf{Holter monitoring (24 to 48 hours):} Captures dynamic changes in the QT interval, T-wave alternans (beat-to-beat variation in T-wave morphology), and silent ventricular arrhythmias during daily activities and sleep.
    \item \textbf{Genetic testing:} Identifies pathogenic mutations in up to 75\% to 80\% of patients with a clinical diagnosis of LQTS, confirming the genotype and enabling cascade screening of family members.
    \item \textbf{Echocardiography:} Performed to rule out structural heart disease, ensuring that the QT prolongation is not secondary to cardiomyopathy or valvular dysfunction.
\end{itemize}

\section*{Management}
The management of LQTS aims to prevent life-threatening arrhythmias through lifestyle modifications, pharmacological therapy, and device implantation.
\begin{itemize}
    \item \textbf{Lifestyle modifications:} Avoiding competitive sports (especially swimming for LQT1), avoiding sudden loud noises (for LQT2), and strictly avoiding QT-prolonging drugs listed on databases such as CredibleMeds.org.
    \item \textbf{Beta-blockers:} The cornerstone of pharmacological therapy for LQT1 and LQT2, with non-selective beta-blockers (specifically nadolol or propranolol) being highly effective in blunting sympathetic triggers and reducing cardiac events.
    \item \textbf{Sodium channel blockers:} Adjunctive therapy for LQT3 (such as mexiletine or flecainide), which blocks the late inward sodium current and directly shortens the QTc interval.
    \item \textbf{Implantable cardioverter-defibrillator (ICD):} Indicated for high-risk patients, including those who have survived a cardiac arrest, or those who experience recurrent syncope despite optimal beta-blocker therapy.
    \item \textbf{Left cardiac sympathetic denervation (LCSD):} A surgical procedure involving the removal of the lower half of the left stellate ganglion, reserved for patients who are refractory to beta-blockers or experience frequent ICD shocks.
    \item \textbf{Electrolyte replacement:} Maintaining high-normal serum potassium and magnesium levels, particularly during gastrointestinal illnesses or when taking diuretics.
\end{itemize}

\section*{Complications}
\begin{itemize}
    \item \textbf{Torsades de pointes:} Polymorphic ventricular tachycardia characterised by a twisting of the QRS complexes around the isoelectric line, which can cause syncope or degenerate into ventricular fibrillation.
    \item \textbf{Sudden cardiac death (SCD):} The most devastating complication, often occurring in young, active individuals who were previously undiagnosed.
    \item \textbf{Implanted device complications:} Lead failure, infection, or inappropriate shocks from an ICD, which can cause significant psychological distress and anxiety.
    \item \textbf{Psychological impact:} High levels of anxiety, depression, and restriction of daily activities due to the fear of arrhythmias or sudden death, both in patients and their families.
    \item \textbf{Adverse effects of beta-blockers:} Fatigue, bradycardia, hypotension, depression, and exercise intolerance, which can affect treatment adherence.
\end{itemize}

\section*{Prognosis and Outcomes}
With early diagnosis and appropriate management, the prognosis for patients with LQTS is generally excellent, with a cardiac event rate of less than 1\% per year in treated patients. Compliance with beta-blockers and lifestyle restrictions is the key determinant of long-term survival. Untreated patients, particularly those with a QTc $>500$~ms or a history of syncope, have a significantly higher risk of sudden cardiac death, emphasising the importance of screening and early intervention.

\section*{Prevention and Self-Care}
\begin{itemize}
    \item \textbf{Avoid QT-prolonging medications:} Regularly check all prescribed and over-the-counter medications against up-to-date lists of QT-prolonging substances.
    \item \textbf{Maintain electrolyte balance:} Stay hydrated and seek medical treatment for vomiting or diarrhoea to prevent low potassium and magnesium levels.
    \item \textbf{Safe exercise habits:} Discuss safe physical activity limits with your cardiologist and ensure you never swim alone if diagnosed with LQT1.
    \item \textbf{Manage fevers promptly:} Treat fevers aggressively with paracetamol, as fever can trigger arrhythmias in certain genotypes (specifically LQT2 and LQT3).
    \item \textbf{Ensure family screening:} Urge first-degree relatives to undergo ECG and genetic evaluation if a mutation is identified.
    \item \textbf{Install a home AED:} Consider purchasing an automated external defibrillator (AED) for the home, and ensure family members are trained in CPR.
\end{itemize}

\section*{Sources and Further Reading}
The following reputable organisations provide further information, drug lists, and clinical guidelines on Long QT syndrome:
\begin{itemize}
    \item Cardiac Society of Australia and New Zealand (CSANZ). \textit{Guidelines for the diagnosis and management of inherited primary arrhythmia syndromes.} https://www.csanz.edu.au/
    \item CredibleMeds. \textit{List of QT-prolonging drugs and avoiding QT prolongation.} https://www.crediblemeds.org/
    \item Heart Foundation Australia. \textit{Understanding sudden cardiac arrest and inherited heart conditions.} https://www.heartfoundation.org.au/
    \item Mayo Clinic. \textit{Long QT syndrome diagnosis and care.} https://www.mayoclinic.org/
    \item SADS Foundation (Sudden Arrhythmic Death Syndromes). \textit{Patient education and advocacy for Long QT syndrome.} https://sads.org/
\end{itemize}
